\documentclass[12pt]{article}

\usepackage[dvips,letterpaper,margin=0.75in,bottom=0.5in]{geometry}
\usepackage{cite}
\usepackage{slashed}
\usepackage{graphicx}
\usepackage{amsmath}
\usepackage{amssymb}
\usepackage{braket}
\usepackage{pythonhighlight}
\usepackage{mathtools}
\begin{document}

\newcommand{\ihbar}{\ensuremath{i \hbar}}
\newcommand{\Pss}{\ensuremath{\Psi^*}}
\newcommand{\dPsidt}{\ensuremath{ \frac{\partial \Psi}{\partial t} }}
\newcommand{\dPsidx}{\ensuremath{ \frac{\partial \Psi}{\partial x} }}
\newcommand{\ddPsidx}{\ensuremath{ \frac{\partial^2 \Psi}{\partial x^2} }}
\newcommand{\dPssdt}{\ensuremath{ \frac{\partial \Psi^*}{\partial t} }}
\newcommand{\dPssdx}{\ensuremath{ \frac{\partial \Psi^*}{\partial x} }}
\newcommand{\ddPssdx}{\ensuremath{ \frac{\partial^2 \Psi^*}{\partial x^2} }}

\newcommand{\dphidt}{\ensuremath{ \frac{d \phi}{dt} }}
\newcommand{\dpsidx}{\ensuremath{ \frac{d \psi}{dx} }}
\newcommand{\ddpsidx}{\ensuremath{ \frac{d^2 \psi}{dx^2} }}


\title{PHY 115L \\ Woods-Saxon Model for the Nucleus \\ using Variational Methods }
\author{Michael Mulhearn}

\maketitle

\section{Introduction}

In this assignment we will develop a numerical model for the atomic
nucleus, based on the Woods-Saxon potential, an effective potential
for the nuclear strong force which binds the nucleus together.  A good
reference for our approach is here: https://arxiv.org/pdf/0709.3525.

\section{System of Units}

We will be looking at neutrons and protons (each with a different
mass) and using their reduced mass (context dependent).  Our system of
units from previous assignments includes the mass of the particle in
the definition of the energy scale, which is now quite unworkable for
us.

We will base our new system of units around two quantities:
$$ \hbar c   = 197.326~\rm MeV \; fm$$
and the fine structure constant:
$$\alpha = \frac{e^2}{4\pi\epsilon_0 \hbar c} = \frac{1}{137.035}$$
Atomic nuclei range in size from $1.6~\rm fm$ to $15~\rm fm$.  Therefore,
we will set our length scale to:
$$L \equiv 1~\rm fm$$
We will set our energy scale to:
$$\epsilon = \frac{\hbar c}{L} = 197.326~\rm MeV$$

The Hamiltonian operator:
\begin{equation}
\hat{H} = -\frac{\hbar^2}{2 M} \; \frac{d^2}{dx^2} \, + \, V(x)
\end{equation}
becomes:
\begin{equation}
\frac{\hat{H}}{\epsilon} = -\frac{L^2}{2 \mu} \; \frac{d^2}{dx^2} \, + \, \frac{V(x)}{\epsilon}
\end{equation}
where $\mu$ is dimensionless quantity related to the mass $M$ by:
$$\mu = \frac{Mc^2}{\epsilon}$$

The masses of the nucleons (neutrons and protons) are:
\begin{eqnarray*}
m_{\rm p} &=& 938.272~\rm MeV/c^2\\
m_{\rm n} &=& 939.565~\rm MeV/c^2\\
\end{eqnarray*}
so in our units:
\begin{eqnarray*}
\mu_{\rm p} &=& 4.75493\\
\mu_{\rm n} &=& 4.76148\\
\end{eqnarray*}

As before, we will be considering {\bf non-periodic} wave functions.
To use the Fourier series approach in this case, we take the period
$a$ to be large compared to $L$, and use the fact that the wave
function approaches zero to approximate a periodic function.


The Fourier series is an expansion in planes waves with momentum $p = \hbar k$ and wavenumbers:
$$k = \frac{2 \pi m}{a}$$
And so the kinetic energy $K(m)$ of a particle described by the plane wave for coefficient $c_m$ is:
$$ K = \frac{p^2}{2M} = \frac{(\hbar k)^2}{2M}$$
and so in our units:
$$ \frac{K}{\epsilon} = \frac{1}{2 \mu} k^2 L^2 = \left(\frac{1}{2\mu}\right) \; (2 \pi m)^2 \; \left( \frac{L}{a}\right)^2$$
This differs from our previous exercises by the factor:
$$\left(\frac{1}{2\mu}\right)$$
which you will have to add to your code.

Using the method of separation of variables (e.g. Griffiths Chapter 4) we write the total wave function as:
$$\Psi_{nlm}(r,\theta,\phi) = \frac{u_{nl}(r)}{r} Y_{lm}(\theta,\phi)$$
for quantum numbers $n,l,m$.  The radial part satisfies 
\begin{equation*}
  \frac{\hat{E}}{\epsilon} u(r) = -\frac{L}{2\mu}^2\frac{d^2 u}{dr^2} + \left(V(r)+ \frac{1}{2\mu}l(l+1) \frac{L^2}{r^2}\right)
\end{equation*}
Notice that this looks exactly like a 1-D particle, but with an effective potential:
$$\frac{V_{\rm eff}(r)}{\epsilon} = V(r) + \frac{1}{2\mu}l(l+1) \frac{L^2}{r^2}$$
The term proportional to $l(l+1)$ is called the centrifugal term, and
it accounts for the additional kinetic energy in the non-radial
direction, due to the angular momentum of the particle.  Not suprisingly (since it is due to kinetic energy) it differs from our previous exercise by the same factor
$$\left(\frac{1}{2\mu}\right)$$
which you will have to add to your code as well.

\section{Woods-Saxon Potential}

We'll use the Woods-Saxon potential:
\begin{equation}
V(r) = -\frac{V_0}{1 + \exp\left( \frac{r-R}{a} \right)}
\end{equation}
where:
\begin{equation}
R = r_0 A^{1/3} 
\end{equation}
and $A$ is the mass number.  Typical values for the parameters are:
\begin{eqnarray*}
V_0 &=& 50 ~ {\rm MeV} \\
r_0 &=& 1.25 ~ {\rm fm}\\
a   &=& 0.5 ~ {\rm fm}\\
\end{eqnarray*}
In our system of units, this becomes:
\begin{equation}
\frac{V(r/L)}{\epsilon} = -\frac{V_0/\epsilon}{1 + \exp\left( \frac{r/L-R/L}{a/L} \right)}
\end{equation}
where
$$R/L = (r_0/L) A^{1/3} $$
and
\begin{eqnarray*}
V_0/\epsilon &=& 0.25 \\
r_0/L &=& 1.25 \\
a/L   &=& 0.5 \\
\end{eqnarray*}

\section{Coulomb Potential}
We will model the charge distribution in the nucleus as a uniform sphere of radius $R$ (with $R$ as above).  So this adds an additional term to the potential:
$$V(r) = Z\frac{e^2}{4\pi\epsilon_0} \frac{F(r/R)}{R}$$
with dimensionless function:
$$F(x) \equiv  \begin{cases}
\frac{3-x^2}{2} & r \leq 1 \\
  1/x & r > 1 \\
\end{cases}
$$
In our system of units this becomes:
$$V(r) = Z \; \alpha \; \left( \frac{L}{R} \right) \; F(r/R)$$



\section{Homework Problems}

\noindent
{\bf Please submit one PDF file (including output) and one python
  notebook file (.ipynb) of your completed assignment.}  If using
C/C++, submit text or PDF of your output and the source code.\\[3pt]


\noindent
{\bf Problem 1:} Building on your results from HW2, find the second
and third excited state of the Harmonic oscillator potential, and compare
the energy your determine for each with the expected value.\\[5pt]

\end{document}
